\documentclass[11pt]{article}

\usepackage[T1]{fontenc}
\usepackage[utf8]{inputenc}
\usepackage{amsmath,amssymb,amsthm}
\usepackage{booktabs}
\usepackage{geometry}
\usepackage{hyperref}
\usepackage{microtype}
\usepackage{xcolor}

\geometry{margin=1in}
\setlength{\emergencystretch}{2em}
\hypersetup{
  colorlinks=true,
  linkcolor=blue!55!black,
  citecolor=blue!55!black,
  urlcolor=blue!55!black
}

\newtheorem{theorem}{Theorem}[section]
\newtheorem{lemma}[theorem]{Lemma}
\newtheorem{proposition}[theorem]{Proposition}
\newtheorem{corollary}[theorem]{Corollary}
\theoremstyle{definition}
\newtheorem{definition}[theorem]{Definition}
\theoremstyle{remark}
\newtheorem{remark}[theorem]{Remark}

\newcommand{\Dfive}{D_5}
\newcommand{\FiveCDC}{\textup{FiveCDC}}
\newcommand{\xor}{\mathbin{\oplus}}
\numberwithin{equation}{section}

\title{A Full Five-Cycle-Double-Cover Boundary Relation\\
for a 288-Vertex Six-Pole}
\author{Atharva Vaidya\\
\small AI-assisted research draft for independent human verification}
\date{July 29, 2026}

\begin{document}
\maketitle

\begin{center}
\fcolorbox{red!65!black}{red!4}{
\parbox{0.91\textwidth}{
\textbf{Resolution status, novelty status, and trust boundary.}
The standard Five-Cycle Double Cover Conjecture remains open.  This note
proves a computer-assisted boundary lemma for one concrete six-pole and
deduces, by a short human argument, a positive theorem for every trivalent
network of copies of that pole.  It is neither a proof nor a disproof of
the general conjecture, and it does not concern the orientable variant.

The novelty claim is provisional.  Searches made before release did not
locate this exact \(D_5\) boundary relation, but the literature contains
substantial adjacent work on FiveCDC superpositions and multipoles.
Specialist prior-art review may show that the network theorem is known or
subsumed by a stronger result.

OpenAI Codex agents, under Atharva Vaidya's direction, selected the
experiment, wrote the programs, discovered and proved the network
consequence, checked the artifacts, and drafted this note.  The independent
checker is implementation-independent, not independent of AI.  No human
peer review is claimed.  A human author must check the proof and provenance,
rerun or reimplement the finite audit, review prior art, and accept
scholarly responsibility before journal or arXiv submission.
}}
\end{center}

\begin{abstract}
Let
\[
        \Dfive=\binom{\{0,1,2,3,4\}}2
\]
be the ten duads, represented by their five-bit characteristic vectors.
A \(\Dfive\)-labelling of a cubic graph is a five-cycle double cover
exactly when the xor of the three incident labels is zero at every vertex.
We compute the complete six-port boundary relation of a concrete
288-vertex \((3,3)\)-pole retained from the \(g=10\) construction
of Karab\'a\v{s}, M\'a\v{c}ajov\'a, Nedela, and \v{S}koviera.

The relation is maximal subject to the global parity obstruction: an
ordered boundary word \((b_0,\ldots,b_5)\in\Dfive^6\) extends through the
pole if and only if \(b_0\xor\cdots\xor b_5=0\).  The positive direction
is certified by 571 explicit internal labellings, one for every orbit of
the 62,560 admissible boundary words under coordinate permutations.
An independently written standard-library checker exhausts the
one-million-word universe and directly checks every certificate; no
negative SAT conclusion is used.

The full relation has a simple all-size consequence.  In every finite
network obtained by attaching all ports of any number of copies of the
pole to trivalent junctions, the connector--junction incidence multigraph
is cubic bipartite.  Three perfect matchings give compatible triangle-duad
labels at every junction and a zero-xor boundary word at every pole.
Consequently every such network has a five-cycle double cover.
\end{abstract}

\section{The exact claim and its scope}

A \emph{cycle} here means an Eulerian edge-subset: every vertex has even
degree in it.  It may be empty or disconnected.  A \emph{five-cycle double
cover} is an indexed family \(C_0,\ldots,C_4\) of such subsets in which
every edge occurs in exactly two members.  Allowing empty members makes
``five'' mean ``at most five.''  Parallel edges, when present, are distinct
edge objects.  The pole used below is loopless; the network construction
may create repeated macro-incidences but no incidence loops.

The standard conjecture, attributed to Celmins and Preissmann, says that
every finite bridgeless graph has such a cover \cite{LiuEtAl2023}.  It
remains open.  In particular, the theorem in this note only eliminates a
concrete superposition family from a counterexample search.

Karab\'a\v{s} et al. construct cyclically 5-edge-connected snarks of
prescribed girth and oddness two by superposition
\cite[Theorem~5.1]{KarabasEtAl2022}.  Our retained \(F_{10}\) pole is the
concrete \(g=10\) instance built with the 70-vertex Harries \((3,10)\)-cage:
the construction first forms their graph \(K_M\), then deletes the two
specified trivalent vertices \(u_1,u_3\), retaining their incident edges
as two ordered three-port connectors.  The resulting pole has 288 vertices,
429 internal edges, and six semiedges.

For reproducibility, this note treats the exact topology in
\texttt{f10-atom.txt} as the formal definition.  Its SHA-256 digest is
\begin{center}
\small\texttt{484503e6318f419457e5160e9f3f7605002621661e82456d57f200c8b6750742}.
\end{center}
The first line records \(288\), \(429\), and \(6\); the second line gives
the six port vertices, grouped as \((0,2,3)\) and \((13,15,16)\); the
remaining lines are the internal edges.  The package records the
construction function from which this frozen topology was exported; the
certificate theorem itself depends only on the frozen topology.

\section{Duad labellings and FiveCDC}

Represent each \(a\in\Dfive\) by its characteristic vector in
\(\mathbb F_2^5\).  Thus every \(a\) has Hamming weight two.

\begin{lemma}[duad equivalence]
\label{lem:duad}
Let \(G\) be a finite loopless graph.  The following data are equivalent:
\begin{enumerate}
\item five Eulerian edge-subsets \(C_0,\ldots,C_4\) covering every edge
      exactly twice;
\item a map \(\lambda:E(G)\to\Dfive\) such that
      \[
          \mathop{\xor}_{e\ni v}\lambda(e)=0
          \qquad\text{for every }v\in V(G).
          \tag{2.1}
      \]
\end{enumerate}
\end{lemma}

\begin{proof}
Given the five subsets, label an edge by the two indices of the subsets
which contain it.  The \(i\)-th coordinate of (2.1) says exactly that
\(C_i\) has even degree at \(v\).  Conversely, from a duad labelling put
\(C_i=\{e:i\in\lambda(e)\}\).  Weight two gives exact double coverage,
and the coordinate equations give the Eulerian condition.
\end{proof}

For a cubic vertex, (2.1) says that the three incident duads form a
coordinate triangle.  For example,
\[
                 01\xor 02\xor 12=0.             \tag{2.2}
\]
For a pole, a boundary word records the labels assigned to its semiedges.
It is \emph{extendible} if labels can be assigned to all internal edges so
that (2.1) holds at every pole vertex, counting its incident semiedge.

\section{The complete boundary relation}

Write \(P\) for the frozen 288-vertex pole and
\[
 R_P=\{(b_0,\ldots,b_5)\in\Dfive^6:
              (b_0,\ldots,b_5)\text{ is extendible through }P\}.
\]

\begin{theorem}[full \(F_{10}\) boundary relation]
\label{thm:boundary}
\[
 R_P=\{(b_0,\ldots,b_5)\in\Dfive^6:
                 b_0\xor\cdots\xor b_5=0\}.       \tag{3.1}
\]
\end{theorem}

\begin{proof}
For necessity, xor (2.1) over all 288 pole vertices.  Every internal edge
appears twice and cancels, leaving precisely the six boundary labels.

The converse is a finite, positive certificate.  There are \(10^6\)
ordered words in \(\Dfive^6\), of which 62,560 have total xor zero.
The symmetric group \(S_5\) acts simultaneously on the five coordinates
of all six duads.  Direct canonicalization partitions the 62,560 admissible
words into 571 orbits.

The frozen file \texttt{s5-representatives.json} contains one canonical
word from each orbit.  The corresponding row of
\texttt{boundary-witnesses.jsonl} contains 435 duads: first the labels of
the 429 internal edges in frozen order, then the six boundary labels.  For
each of the 571 rows, the checker verifies that every entry is a duad, that
the claimed boundary is the orbit representative, and that the three
labels incident with each of the 288 vertices xor to zero.  A permutation
of the five coordinates carries a verified witness to a witness for every
other word in the same orbit.  Hence every word satisfying (3.1) extends.
\end{proof}

\begin{remark}[what is and is not trusted]
The proof uses no solver's assertion of unsatisfiability.  CaDiCaL was used
only as a producer of positive labellings.  The checker is an independently
written Python program using only the standard library.  It reconstructs
the entire one-million-word universe, recomputes the 571 orbits, and
directly checks 248,385 edge labels and 164,448 vertex equations.  A full
replay also recompiles and reruns the producer.  The witness file has
SHA-256 digest
\[
\small\texttt{3e7328b4d02a459f6d1f85b8f357aa5ee70c38d81138bcea5b8ecaef84b2961d}.
\]
\end{remark}

\begin{corollary}[emergent port symmetry]
The boundary relation \(R_P\) is invariant under the full symmetric group
\(S_6\) on the ports.
\end{corollary}

\begin{proof}
Condition (3.1) is invariant under every permutation of its six entries.
\end{proof}

This is symmetry of the relation, not a claim that every port permutation
is induced by an automorphism of the underlying pole.  Quotienting the
relation by \(S_5\times S_6\) leaves 11 orbits.

\section{Every trivalent network is positive}

\begin{definition}
An \emph{\(F_{10}\) trivalent network} is obtained from finitely many
disjoint copies of \(P\) by attaching every port to exactly one new
trivalent junction, with exactly three port edges incident with every
junction.  A junction may receive several ports from the same pole or
from the same connector triple.
\end{definition}

\begin{theorem}[universal network theorem]
\label{thm:network}
Every finite \(F_{10}\) trivalent network has a five-cycle double cover.
\end{theorem}

\begin{proof}
Form a bipartite incidence multigraph \(H\).  Its left vertices are the
two connector triples of every copy of \(P\); its right vertices are the
trivalent junctions; and its edges are the individual port-to-junction
incidences.  Every left vertex has degree three because a connector has
three ports, and every right vertex has degree three by construction.
Repeated incidences give parallel edges, which are retained separately.
Thus \(H\) is a cubic bipartite multigraph.

Every cubic bipartite multigraph is 3-edge-colourable.  Indeed, for a set
\(S\) of vertices in one part, its \(3|S|\) incident edges end in
\(N(S)\), which can receive at most \(3|N(S)|\) of them.  Hall's condition
holds, so \(H\) has a perfect matching.  Removing it leaves a 2-regular
bipartite multigraph, whose even cycles split into two further perfect
matchings.  Call the three colours \(a,b,c\).

Map these colours to the triangle duads
\[
                  q_a=01,\qquad q_b=02,\qquad q_c=12.
\]
At every junction each colour occurs once, so (2.2) gives the vertex
condition.  At every connector each colour also occurs once.  The two
connectors belonging to one pole therefore contain each of
\(q_a,q_b,q_c\) exactly twice across their six ports, and their total xor
is zero.  By Theorem~\ref{thm:boundary}, this boundary word extends through
that pole.  Choose an extension independently for every pole.  The
resulting duad labelling satisfies (2.1) at every vertex of the expanded
network.  Lemma~\ref{lem:duad} converts it to a five-cycle double cover.
\end{proof}

\begin{remark}
The positive labelling does not require the expanded graph to be simple,
connected, bridgeless, or high-girth.  Those restrictions matter only when
the network is proposed as a candidate counterexample.
\end{remark}

\section{Reproducibility}

The complete package is publicly available in the draft repository and is
also located, relative to this source tree, in
\begin{center}
\texttt{scratch/f10-boundary-relation-20260729/}.
\end{center}
The public review location is
\href{https://github.com/AtharvaVaidya/mathematics-research/pull/5}
{\texttt{github.com/AtharvaVaidya/mathematics-research/pull/5}}.
From that directory, the fast independent audit is
\begin{verbatim}
shasum -a 256 -c SHA256SUMS
python3 verify_boundary_relation.py \
  f10-atom.txt s5-representatives.json \
  boundary-witnesses.jsonl relation-summary.json
\end{verbatim}
On the frozen package it reports
\begin{verbatim}
ordered_duad_words=1000000 xor_zero=62560 S5_orbits=571
positive=571 negative=0
directly checked 571 witnesses, 248385 edge labels,
and 164448 vertex xor equations
\end{verbatim}
The package ledger itself has SHA-256 digest
\[
\small\texttt{5b3abc755f6879dc4c938d5da7f2ae507e47457cc735fe80c9ea58a4226ed6f6}.
\]

The command \texttt{python3 verify.py} performs an end-to-end replay.  It
first invokes the independent checker, then regenerates the orbit census,
compiles the C++20 producer, reruns all 571 positive SAT queries, and checks
the newly produced witnesses directly.  This replay currently expects the
CaDiCaL headers and static library in the Homebrew paths recorded by the
script.  The direct checker has no non-standard dependency.

\section{Relationship to prior work and remaining questions}

The pole topology comes from Karab\'a\v{s} et al.
\cite{KarabasEtAl2022}.  Their theorem concerns girth, oddness, colouring
defect, and proper superposition; it does not state the exact
\(\Dfive\)-boundary relation proved here.  More broadly, Liu, Hao, Luo,
and Zhang give sufficient conditions for superpositions to have a
five-cycle double cover and verify several snark families
\cite{LiuEtAl2023}.  M\'a\v{c}ajov\'a and \v{S}koviera develop the modern
formalism of proper multipoles and superposition
\cite{MacajovaSkoviera2021}.  These are the closest sources found in the
initial search and must be compared in depth before asserting novelty.

The contribution claimed by this draft is therefore deliberately narrow
and provisional:
\begin{enumerate}
\item the explicit complete \(\Dfive\) boundary relation of one frozen
      288-vertex six-pole;
\item independently checkable positive witnesses for all boundary orbits;
\item the short cubic-bipartite incidence proof covering every trivalent
      network of copies of that pole.
\end{enumerate}
The most natural next question is whether the same full boundary relation
holds for a structural class of the \(F_g\) poles, rather than only this
concrete \(g=10\) instance.  A negative answer would also be useful:
non-full boundary relations are potential building blocks for
certificate-producing counterexample searches.

\section*{Acknowledgment and AI-use disclosure}

This manuscript is an openly AI-assisted research artifact.  OpenAI Codex
agents made material contributions to the mathematical search, software,
proof derivation, verification design, literature search, and prose.  The
repository history and frozen certificates are provided so that the claims
can be assessed without trusting the AI system.  The named human has not,
by publication of this draft alone, represented that the work has passed
independent expert review.

\begin{thebibliography}{9}

\bibitem{KarabasEtAl2022}
J.~Karab\'a\v{s}, E.~M\'a\v{c}ajov\'a, R.~Nedela, and M.~\v{S}koviera,
\emph{Girth, oddness, and colouring defect of snarks},
Discrete Mathematics 345 (2022), Article 113040.
\href{https://doi.org/10.1016/j.disc.2022.113040}
{doi:10.1016/j.disc.2022.113040}.

\bibitem{LiuEtAl2023}
S.~Liu, R.-X.~Hao, R.~Luo, and C.-Q.~Zhang,
\emph{5-Cycle Double Covers, 4-Flows, and Catlin Reduction},
SIAM Journal on Discrete Mathematics 37 (2023), 253--267.
\href{https://doi.org/10.1137/22M1472425}
{doi:10.1137/22M1472425}.

\bibitem{MacajovaSkoviera2021}
E.~M\'a\v{c}ajov\'a and M.~\v{S}koviera,
\emph{Superposition of snarks revisited},
European Journal of Combinatorics 91 (2021), Article 103220.
\href{https://doi.org/10.1016/j.ejc.2020.103220}
{doi:10.1016/j.ejc.2020.103220}.

\end{thebibliography}

\end{document}
